\subsubsection{Quaternion Basis Under Star and H}
Hamilton's quaternion units are also reproduced by the same sigma-number-two
basis used in Eq.\ \eqref{eq:exp-bivector}, so the rotation algebra can be
written in quaternion form without changing any multiplication rule
\cite{hamilton1844}:
\begin{equation}
\begin{aligned}
\pv{q}_1&=-\pv{\sigma}_2\pv{\sigma}_3,\qquad
\pv{q}_2=-\pv{\sigma}_3\pv{\sigma}_1,\\
\pv{q}_3&=-\pv{\sigma}_1\pv{\sigma}_2,\qquad
\vct{q}:=\vcol{\pv{q}_1,\pv{q}_2,\pv{q}_3}=-i\sigv.
\end{aligned}
\label{eq:quaternion-basis}
\end{equation}
The standard quaternion identities are satisfied by these basis elements:
\begin{equation*}
\pv{q}_1^2=\pv{q}_2^2=\pv{q}_3^2=\pv{q}_1\pv{q}_2\pv{q}_3=-1.
\end{equation*}
Thus the quaternion units are not new generators; the existing
sigma-number-two basis is only relabeled.

If \(\vct{B}=\vcol{B_1,B_2,B_3}\in\R^3\), then
\(\vct{B}\cdot\vct{q}=-i\,\vct{B}\cdot\sigv\). Since
\(\pv{q}_k=-i\pv{\sigma}_k\), the following relations are obtained
immediately from the previously defined star and \(H\) rules:
\begin{equation*}
\vct{q}^*=\vct{q},
\qquad
\vct{q}^{\mathsf H}=-\vct{q}.
\end{equation*}
Thus the quaternion basis is fixed by the conjugate operator, while its
sign is changed by the Hermitian operation.

\subsubsection{Quaternion Products}
For \(a\in\R\) and \(\vct{B}\in\R^3\), the quaternion is defined by
\begin{equation}
\pv{Q}:=a+\vct{B}\cdot\vct{q}
=a-i\,\vct{B}\cdot\sigv.
\label{eq:quaternion-paravector}
\end{equation}
A real scalar coefficient and a purely imaginary vector coefficient are
carried, so \(\pv{Q}\in\Pspace(\C)\).
In complex-paravector notation, the same quaternion may be written as
\begin{equation*}
\pv{Q}=\pv{X}-i\pv{Y},
\qquad
\pv{X}=a+\vct{A}\cdot\sigv,
\qquad
\pv{Y}=b+\vct{B}\cdot\sigv,
\end{equation*}
with the quaternion specialization
\begin{equation*}
\vct{A}=0,
\qquad
b=0.
\end{equation*}
Under this specialization, \(\pv{X}=a\) is purely scalar and
\(\pv{Y}=\vct{B}\cdot\sigv\) is purely vector. Because
\(\vct{q}^*=\vct{q}\),
\begin{equation*}
\pv{Q}^*=\pv{Q}.
\end{equation*}
Therefore
\begin{align}
\pv{Q}\pv{Q}^*&=\pv{Q}^2\notag\\
&=(a-i\,\vct{B}\cdot\sigv)^2\notag\\
&=a^2-2ia\,\vct{B}\cdot\sigv+(-i)^2(\vct{B}\cdot\sigv)^2\notag\\
&=a^2-\vct{B}^2-2ia\,\vct{B}\cdot\sigv.
\label{eq:quaternion-square}
\end{align}
The factor \(a\) must be carried by the linear term. This is just the
sigma-number-two multiplication law written in quaternion language.

Since \(\vct{q}^{\mathsf H}=-\vct{q}\), the Hermitian counterpart of \(\pv{Q}\) is
\begin{equation*}
\pv{Q}^{\mathsf H}=a-\vct{B}\cdot\vct{q}
=a+i\,\vct{B}\cdot\sigv.
\end{equation*}
The result of multiplication by the Hermitian counterpart is
\begin{equation}
\pv{Q}\pv{Q}^{\mathsf H}=a^2+\vct{B}^2.
\label{eq:quaternion-norm}
\end{equation}
Thus the standard positive quaternion norm square is recovered by the
Hermitian product. Here ``norm square'' means the nonnegative real number
that measures the coefficient magnitude. The ordinary algebra square is
instead recorded by
Eq.\ \eqref{eq:quaternion-square}.

\subsubsection{Quaternion Form of the Rotor Law}
A quaternion satisfying \(\pv{Q}\pv{Q}^{\mathsf H}=1\) is called a unit
quaternion. For such a quaternion and a real vector
\(\vct{r}\in\R^3\), the rotation law is
\begin{equation}
\vct{r}'\cdot\vct{q}
=\pv{Q}(\vct{r}\cdot\vct{q})\pv{Q}^{\mathsf H}.
\label{eq:quaternion-rotation}
\end{equation}
This is the quaternion version of the rotor action
\(\vct{r}'\cdot\sigv=\herm{\pv{R}}(\vct{r}\cdot\sigv)\pv{R}\). With the
choice
\begin{equation*}
\pv{Q}=\cos\!\left(\frac{\theta}{2}\right)
 +\sin\!\left(\frac{\theta}{2}\right)\vct{u}\cdot\vct{q}
\end{equation*}
the relation \(\pv{Q}=\pv{R}^{\mathsf H}\) is obtained because
\(\vct{q}=-i\sigv\). Consequently
\[
\pv{Q}(\vct{r}\cdot\vct{q})\pv{Q}^{\mathsf H}
=\pv{R}^{\mathsf H}(\vct{r}\cdot\vct{q})\pv{R},
\]
which has exactly the same orientation as the sigma-vector rotor law. The
equivalent two-sided quaternion product
\(\pv{Q}^{\mathsf H}(\vct{r}\cdot\vct{q})\pv{Q}\) is required under the
alternative choice \(\pv{Q}=\pv{R}\). Thus quaternion rotations are
the same circular half-angle algebra, with the order fixed by the selected
choice of positive rotation direction.
